\subsection{Requirements for Reusability and Standardization}
\label{subsec:wiederverwendbarkeit-standardisierung}

The baseline is intended to define shared structures and public workflows while
allowing controlled profile variants. Centralized maintenance prevents copies
from diverging. Table~\ref{tab:anforderungen-template} summarizes the verifiable
requirements.

\begin{table}[H]
  \centering
  \scriptsize
  \renewcommand{\arraystretch}{1.08}
  \caption{Key requirements for the technology template}
  \label{tab:anforderungen-template}
  \begin{tabularx}{\textwidth}{@{}p{0.08\textwidth}p{0.22\textwidth}p{0.39\textwidth}>{\raggedright\arraybackslash}X@{}}
    \toprule
    \textbf{ID} & \textbf{Requirement} & \textbf{Operationalized objective} & \textbf{Basis for subsequent verification} \\
    \midrule
    A-01 & Baseline & It must be possible to generate the basic structure. & Project generation \\
    A-02 & Structure & Responsibilities are assigned to consistent locations. & Repository structure \\
    A-03 & Profiles & Only designated components are activated. & Profile model \\
    A-04 & Workflows & Checks, initialization, installation, startup, testing, and building share common entry points. & Tooling \\
    A-05 & Maintainability & Shared components are maintained centrally. & Repository strategy \\
    A-06 & Traceability & Changes and configuration are versioned and documented. & Version state, documentation \\
    A-07 & Extensibility & Optional capabilities have documented dependencies. & Capability model \\
    \bottomrule
  \end{tabularx}
  \smallskip
  {\footnotesize Source: Author's own compilation.\par}
\end{table}

The stated verification basis is subsequently used to assess these target
criteria.
\beginnerinput{workspace/statement/200_anforderungen/220}
